% GENERATED FILE. Edit canonical lessons, then run npm run generate:books.
% canonical-source-hash: fnv1a64:a2461bf89cfc590f
% canonical-lessons: LA-C20-quid-tibi-nomen, LA-C20-name-case-variation

\chapter{Exchanging Names}
\label{ch:la-exchanging-names}

This chapter is generated from the canonical micro-lessons used by Language Ladder.
Each section stays independently resumable and preserves the authored prerequisite order.

\section[quid tibi nōmen est? / mihi nōmen est...]{quid tibi nōmen est?, mihi nōmen est... — "what's your name?", the dative way}

\label{lesson:LA-C20-quid-tibi-nomen}



\begin{quote}
\textbf{Warm-up.} {[}PAUSE 2s{]} English says you \textbf{have} a name. German says you're \textbf{called} something. Latin does neither — it says the name \textbf{belongs to you}, using nothing but a case ending.
\end{quote}

\subsection*{You'll want to know: quid tibi nōmen est? — literally "what is the name TO YOU?"}
\textbf{Quid tibi nōmen est?} ("\textbf{what's your name?}") is, word for word, "\textbf{what} {[}is{]} \textbf{the-name} {[}to{]} \textbf{you}?" — no verb for "have," no verb for "be called." The name simply \textbf{exists}, and \textbf{tibi} ("to you," dative case) says who it belongs to.

\subsection*{You'll want to know: mihi nōmen est... — the same construction, in reply}
The reply follows the identical pattern: \textbf{mihi nōmen est} {[}name{]} — "{[}the-name{]} {[}is{]} to-me {[}name{]}." This is Latin's genuine \textbf{"dative of possession"}: instead of a possessive verb, the possessor simply sits in the dative case, and the thing possessed — here, \textbf{nōmen}, "name" — stays in the nominative, as the sentence's real subject.

\subsection*{Guided Practice}
{[}PAUSE 1s{]}

\begin{itemize}
\raggedright
  \item {[}YOU SAY: "quid tibi nōmen est?" — literally "what's the name to you?"{]}
  \item {[}YOU SAY: "mihi nōmen est..." — "the name is to me..."{]}
  \item {[}YOU SAY: the case doing the work — dative \emph{tibi}/\emph{mihi}, not a verb{]}
\end{itemize}

\begin{tcolorbox}[breakable,colback=teal!4,colframe=teal!35!black,title={Wrap-up recall}]
{[}PAUSE 3s{]} What case does the possessor sit in, in \textbf{mihi nōmen est}? (\textbf{Dative} — \emph{mihi}, "to me.") Does Latin use a verb like "have" or "be called" for this? (\textbf{No} — the name simply exists; possession is shown by the dative case alone.) What noun remains the sentence's grammatical subject? (\textbf{Nōmen}, "name.")
\end{tcolorbox}
\section[puerō nōmen est Mārcus / Mārcō]{Mārcus or Mārcō? — a real classical variation}

\label{lesson:LA-C20-name-case-variation}



\begin{quote}
\textbf{Warm-up.} {[}PAUSE 2s{]} The possessor in Latin's naming construction is reliably dative. The personal name beside \textbf{nōmen}, however, is not locked to one case in every author.
\end{quote}

\begin{grammarlens}[title={Grammar Lens: The stable frame}]
The dative possessor appears across authors from Plautus through Livy and Sallust:

\begin{quote}
\textbf{puerō nōmen est ...} — to the boy the name is ...
\end{quote}

That \textbf{puerō} does not change. It marks the person who possesses the name.
\end{grammarlens}

\begin{grammarlens}[title={Grammar Lens: Two habits for the personal name}]
\textbf{Cicero} tends to keep the personal name nominative, agreeing with \textbf{nōmen}, the sentence's subject:

\begin{quote}
\textbf{puerō nōmen est Mārcus}
\end{quote}

\textbf{Livy} and \textbf{Sallust} more often pull the personal name into the dative to match the possessor:

\begin{quote}
\textbf{puerō nōmen est Mārcō}
\end{quote}

The underlying dative-of-possession construction is the same. The choice between \textbf{Mārcus} and \textbf{Mārcō} is authorial variation, not evidence that one sentence has stopped using the dative of possession.
\end{grammarlens}

\subsection*{Guided Practice}
{[}PAUSE 1s{]}

\begin{itemize}
\raggedright
  \item {[}YOU SAY: "puerō" — to the boy; always the dative possessor{]}
  \item {[}YOU SAY: "Mārcus" — Cicero's usual nominative pattern{]}
  \item {[}YOU SAY: "Mārcō" — the matching dative found more often in Livy/Sallust{]}
\end{itemize}

\begin{tcolorbox}[breakable,colback=teal!4,colframe=teal!35!black,title={Wrap-up recall}]
{[}PAUSE 3s{]} Which case remains constant for the possessor? (\textbf{Dative.}) Which author tends toward nominative \textbf{Mārcus}? (\textbf{Cicero.}) Which authors more often use matching dative \textbf{Mārcō}? (\textbf{Livy and Sallust.}) Do the two forms represent different possession constructions? (\textbf{No.})
\end{tcolorbox}
